\appendix

\chapter{Status, verification and provenance}\label[appendix]{app:status}

\section{Status vocabulary}
Every theorem, proposition and observation whose proof or evidence is not purely analytic carries its status in its title. \emph{Analytic}: proved in the text from the stated hypotheses and classical results. \emph{Computer-assisted}: the proof reduces to finite statements certified by exact computation. \emph{Finite-exact}: an exact computation on a stated finite domain, with no claim beyond it. \emph{Bounded}: a statement established over a stated bounded range, with no extrapolation. Statements that were found false are kept as such in the text and named \emph{falsified}; nothing is stated as conjectural without being marked so. All computations use exact integer or \texttt{Fraction} arithmetic and linear algebra over $\F$ on bit-packed rows; floating point appears only in drawing figures from exact data.

\section{What was recomputed for this monograph}
The package that accompanies the monograph contains a single verification module and five check scripts (\texttt{verify/gigmini.py}, \texttt{run\_checks.py}--\texttt{run\_checks5.py}) with their recorded outputs. They re-derive the histories of $72$ and Lemma~F, the two anchors and their windows, witnesses W1, W2, W5 and W6, the radius-two identity, instances of the nerve theorem and of Alexander duality on random admissible sets, the adjacency lemmas (F1)--(F3), the triple counterexample, the lift calculus on the anchor, the certificate $(z,\varphi)$ with the families $B_0,\dots,B_3$, the two gate counterexamples, and the nerve and core computations on the named systems and on all $290$ known positives. A stand-alone checker that evaluates the $290$ stored certificates was run for this build, with the result that all checks pass, and two deliberately corrupted certificate files are each rejected; the recorded outputs are in the package. Output files are listed in the package record \texttt{records/VERIFICATION.md}.

\section{Computations that were not repeated}
The following large inherited computations were not re-executed; the numbers quoted in the text are taken from their recorded outputs and are labelled as such where they occur: the six-label census of $346{,}859$ regimes and the seven-label scan of $1{,}560{,}780$ regimes, the three-lift search, the nerve table of \cref{tab:D21} ($219{,}145$ patterns, $151$ positive), the held-out prediction rounds of \cref{obs:atlas-pred}, the atlas rounds of \cref{obs:atlas}, the out-of-sample information values of \cref{obs:oos}, and the exact history counts of $72^r$ for $r=8,9,10$.

\section{Results not proved here}
The following are verified on a finite domain but have no proof for all parameters: face-recognition soundness for all $k$ (\cref{thm:face-recognition}, computer-assisted for $k\le7$); the classification of $6$-cycles through a vertex for $k>8$ (\cref{thm:skeleton}); non-degeneracy of off-diagonal lifts in general and injectivity of $\first_\ast$ (\cref{obs:nondeg}, \cref{thm:JX}(J-7)); injectivity of the history count on the comparable window in general (\cref{obs:anchored}); the connectedness parts of (F1)--(F2) by hand, and the pairwise theorem for $k\ne6$ (\cref{thm:pairwise}). The strong form of gate necessity is false (\cref{prop:gate-nec}); only the exact criterion of \cref{thm:exact-survival} is claimed. Additional corrected or falsified claims in the earlier computations: a claimed gate of $18{,}980$ regimes, the signature $S_A=(N,S_{BB})$ as a separator (\cref{obs:preregistered}), the principle that $(Q,S_{s2})$ separates in general (\cref{obs:S2-falsified}), and a ladder claim for $\beta_1$, now restricted to prefix-defined sets with $k\le8$ (\cref{thm:obs-ladder}). The cellular-resolution reading of \cref{sec:relations} is an interpretation only.

\section{Source releases and manifest hashes}
The mathematics was assembled from earlier frozen releases, listed with the first sixteen hexadecimal digits of the sha256 of each release's manifest file, as present in the build workspace.
\begin{center}\footnotesize\setlength{\tabcolsep}{5pt}
\begin{tabular}{@{}llll@{}}\toprule Release & Hash & Release & Hash\\\midrule
GIG-I v001 & 28385614d8999629 & GIG-INFO v001 & cb03261b27c3ffc9\\
GIG-I v002 & 0212b57ee3306c6f & GIG-INFO v002 & 7709cc7a02d31fa4\\
GIG-I v003 & fe59b68f8dcbb0b8 & GIG-INFO v003 & 6e79f1013cca47bb\\
GIG-I v004 & 26e14d978aafe3ff & GIG-INFO v004 & abf47ee29f59bebe\\
GIG-I v005 & aea157c8fd126a00 & GIG-INFO v005 & c01a7298716615d0\\
GIG-I v005 S2 & 9558b1071737ff35 & GIG-INFO v006 & 9e4c14ede3e65fe6\\
GIG-I v005 S3 & e38b5e73d4869902 & & \\
GIG-I v006 & ac17c9efddd9f692 & & \\
GIG-I v007 & 2d1d3aa0094c4a8c & & \\
\bottomrule\end{tabular}\end{center}
\paragraph{Erratum on parent hashes.} The fifth GIG-INFO release recorded the same parent hash for GIG-I v004 and v005, because a path helper rewrote only the first occurrence on one line. The true GIG-I v004 manifest hash is the one listed above. The sixth release re-tested a one-token repair: all mathematical tests passed and the affected output regenerated identically, and no result depends on a module of GIG-I v004 alone. Of the $14$ parent manifests recorded in the sixth release, $11$ were matched to the workspace copies; $3$ did not match or could not be matched and remain unresolved.
\paragraph{Repository links.} The fifteen research releases on which this monograph is based are published as a single collection at a fixed commit~\cite{UNNSGIG}; the commit identifies the snapshot, not whatever later appears on the main branch. The present monograph is not part of that commit. Each release is also identified by its manifest hash in the table above. The reference collection on the mathematical foundations of the Substrate is pinned separately and cited as~\cite{UNNSFound}. A statement-by-statement index of the releases in which each result was established, with links at the pinned commit, accompanies the package (\texttt{records/THEOREM\_SOURCE\_INDEX\_v002.md}).

\chapter{Certificates of the anchor hole}\label[appendix]{app:certs}
\section{The anchor certificate}
Let $\lambda_A=(2,3,4,5,8,9)$, with $K_z=288$. A \emph{certificate} for a non-trivial class in $H_2(X;\F)$ is a pair $(z,\varphi)$ with $z$ a $2$-cycle supported in $X$ and $\varphi$ a $2$-cochain with $\varphi\circ\partial_3=0$ on every $3$-face of $P_6$ and $\varphi(z)=1$ (\cref{lem:certificate}). The verification consists of three checks, each independent of the others: $\partial z=0$; $\varphi(\partial F)=0$ for each of the $540$ three-faces $F$ of $P_6$; $\varphi(z)=1$. Because $\varphi$ is defined on $P_6$ and not on $X$, it vanishes on the boundaries of all $3$-faces of any subcomplex: $[z]\ne0$ in every $X$ containing $z$, which is how the survival results are used.

\subsection*{The cycle $z$ ($96$ two-faces)}
Each face is written as an ordered set partition of $\{0,\dots,5\}$ into four blocks (blocks separated by \texttt{|}); block sizes $(3,1,1,1)$ in some order give a hexagon and $(2,2,1,1)$ a square. There are $64$ squares and $32$ hexagons. The list is sorted lexicographically.
\begin{center}\footnotesize\setlength{\tabcolsep}{6pt}
\begin{tabular}{llll}
\texttt{012|3|4|5} & \texttt{014|2|3|5} & \texttt{01|24|3|5} & \texttt{01|2|34|5} \\
\texttt{025|1|3|4} & \texttt{02|15|3|4} & \texttt{02|1|35|4} & \texttt{02|1|3|45} \\
\texttt{045|1|2|3} & \texttt{04|15|2|3} & \texttt{04|1|25|3} & \texttt{04|1|2|35} \\
\texttt{05|12|3|4} & \texttt{05|14|2|3} & \texttt{05|1|24|3} & \texttt{05|1|2|34} \\
\texttt{0|125|3|4} & \texttt{0|12|35|4} & \texttt{0|12|3|45} & \texttt{0|145|2|3} \\
\texttt{0|14|25|3} & \texttt{0|14|2|35} & \texttt{0|15|24|3} & \texttt{0|15|2|34} \\
\texttt{0|1|245|3} & \texttt{0|1|24|35} & \texttt{0|1|25|34} & \texttt{0|1|2|345} \\
\texttt{124|0|3|5} & \texttt{12|04|3|5} & \texttt{12|0|34|5} & \texttt{14|02|3|5} \\
\texttt{1|024|3|5} & \texttt{1|02|34|5} & \texttt{245|0|1|3} & \texttt{24|15|0|3} \\
\texttt{24|1|05|3} & \texttt{24|1|0|35} & \texttt{24|5|01|3} & \texttt{25|04|1|3} \\
\texttt{25|0|14|3} & \texttt{25|0|1|34} & \texttt{2|013|4|5} & \texttt{2|035|1|4} \\
\texttt{2|03|15|4} & \texttt{2|03|1|45} & \texttt{2|05|13|4} & \texttt{2|0|135|4} \\
\texttt{2|0|13|45} & \texttt{2|134|0|5} & \texttt{2|13|04|5} & \texttt{2|14|03|5} \\
\texttt{2|1|034|5} & \texttt{2|345|0|1} & \texttt{2|34|01|5} & \texttt{2|34|05|1} \\
\texttt{2|34|0|15} & \texttt{2|35|04|1} & \texttt{2|35|0|14} & \texttt{2|3|014|5} \\
\texttt{2|3|045|1} & \texttt{2|3|04|15} & \texttt{2|3|05|14} & \texttt{2|3|0|145} \\
\texttt{2|45|03|1} & \texttt{2|45|0|13} & \texttt{2|4|135|0} & \texttt{2|4|13|05} \\
\texttt{2|4|15|03} & \texttt{2|4|1|035} & \texttt{2|4|35|01} & \texttt{2|4|3|015} \\
\texttt{2|4|5|013} & \texttt{2|5|034|1} & \texttt{2|5|03|14} & \texttt{2|5|04|13} \\
\texttt{2|5|0|134} & \texttt{45|02|1|3} & \texttt{45|0|12|3} & \texttt{4|015|2|3} \\
\texttt{4|01|25|3} & \texttt{4|01|2|35} & \texttt{4|125|0|3} & \texttt{4|12|05|3} \\
\texttt{4|12|0|35} & \texttt{4|15|02|3} & \texttt{4|1|025|3} & \texttt{4|1|02|35} \\
\texttt{4|25|01|3} & \texttt{4|5|012|3} & \texttt{5|024|1|3} & \texttt{5|02|14|3} \\
\texttt{5|02|1|34} & \texttt{5|04|12|3} & \texttt{5|0|124|3} & \texttt{5|0|12|34} \\
\end{tabular}\end{center}

\subsection*{The cochain $\varphi$ (six hexagons)}
\texttt{245|0|1|3}, \texttt{245|0|3|1}, \texttt{245|1|0|3}, \texttt{245|1|3|0}, \texttt{245|3|0|1}, \texttt{245|3|1|0}.

\subsection*{The set EXTRA ($20$ prefix sets) and the minimal families}
Prefix sets are written as strings of the labels (ranks) they contain. EXTRA $=\{3, 03, 13, 013, 024, 34, 034, 134, 0134, 15, 125, 35, 035, 135, 0135, 1235, 345, 0345, 1345, 01345\}$.
\begin{itemize}
\item $B_0$ (1 sets): 024.
\item $B_1$ (3 sets): 15, 125, 1235.
\item $B_2$ (15 sets): 3, 03, 13, 013, 34, 034, 134, 0134, 35, 035, 0135, 345, 0345, 1345, 01345.
\item $B_3$ (17 sets): 3, 03, 13, 013, 34, 134, 0134, 15, 35, 035, 135, 0135, 1235, 345, 0345, 1345, 01345.
\end{itemize}

The Boolean statement of \cref{thm:exact-survival} says that $z$ bounds in $X_\lambda(m)$ with $K_z\mid m$ iff the family $\mathfrak F_\lambda(m)\cap\mathrm{EXTRA}$ contains one of $B_0,\dots,B_3$. Each $B_i$ was checked to bound and to be minimal under removal of one set; the criterion agreed with exact filling on $400$ random subfamilies of the $20$ sets. The count $589{,}839$ of the $2^{20}$ families which bound is inherited and was not recounted.

\begin{thebibliography}{99}\addcontentsline{toc}{chapter}{Bibliography}
\small
\bibitem{BarmakMinian} J.~A.~Barmak and E.~G.~Minian, Strong homotopy types, nerves and collapses, \emph{Discrete Comput.\ Geom.} \textbf{47} (2012), 301--328.
\bibitem{BayerSturmfels} D.~Bayer and B.~Sturmfels, Cellular resolutions of monomial modules, preprint, arXiv:alg-geom/9711023. (Journal reference not verified.)
\bibitem{BjornerBrenti} A.~Bj\"orner and F.~Brenti, \emph{Combinatorics of Coxeter Groups}, Graduate Texts in Mathematics \textbf{231}, Springer, 2005.
\bibitem{Bjorner} A.~Bj\"orner, Topological methods, in: \emph{Handbook of Combinatorics}, vol.~II, North-Holland, 1995. (Page range not verified.)
\bibitem{BjornerTancer} A.~Bj\"orner and M.~Tancer, Combinatorial Alexander duality --- a short and elementary proof, \emph{Discrete Comput.\ Geom.} \textbf{42} (2009), 586--593. doi:10.1007/s00454-008-9102-x.
\bibitem{BlindManiLevitska} R.~Blind and P.~Mani-Levitska, Puzzles and polytope isomorphisms, \emph{Aequationes Math.} \textbf{34} (1987), 287--297.
\bibitem{Borsuk} K.~Borsuk, On the imbedding of systems of compacta in simplicial complexes, \emph{Fund.\ Math.} \textbf{35} (1948), 217--234.
\bibitem{Carlsson09} G.~Carlsson, Topology and data, \emph{Bull.\ Amer.\ Math.\ Soc.} \textbf{46} (2009), 255--308.
\bibitem{CarlssonZomorodian} G.~Carlsson and A.~Zomorodian, The theory of multidimensional persistence, \emph{Discrete Comput.\ Geom.} (2009). (Volume and pages not verified.)
\bibitem{CoverThomas} T.~M.~Cover and J.~A.~Thomas, \emph{Elements of Information Theory}, 2nd ed., Wiley, 2006.
\bibitem{EagonReiner} J.~A.~Eagon and V.~Reiner, Resolutions of Stanley--Reisner rings and Alexander duality, \emph{J.\ Pure Appl.\ Algebra} \textbf{130} (1998), 265--275. doi:10.1016/S0022-4049(97)00097-2.
\bibitem{ELZ02} H.~Edelsbrunner, D.~Letscher and A.~Zomorodian, Topological persistence and simplification, \emph{Discrete Comput.\ Geom.} \textbf{28} (2002), 511--533.
\bibitem{GasharovPeevaWelker} V.~Gasharov, I.~Peeva and V.~Welker, The lcm-lattice in monomial resolutions, \emph{Math.\ Res.\ Lett.} \textbf{6} (1999), 521--532.
\bibitem{Hartley28} R.~V.~L.~Hartley, Transmission of information, \emph{Bell System Tech.\ J.} \textbf{7}(3) (1928), 535--563.
\bibitem{Hatcher} A.~Hatcher, \emph{Algebraic Topology}, Cambridge University Press, 2002.
\bibitem{Kalai88} G.~Kalai, A simple way to tell a simple polytope from its graph, \emph{J.\ Combin.\ Theory Ser.\ A} \textbf{49} (1988), 381--383. doi:10.1016/0097-3165(88)90064-7.
\bibitem{LindMarcus} D.~Lind and B.~Marcus, \emph{An Introduction to Symbolic Dynamics and Coding}, Cambridge University Press, 1995.
\bibitem{MillerSturmfels} E.~Miller and B.~Sturmfels, \emph{Combinatorial Commutative Algebra}, Graduate Texts in Mathematics \textbf{227}, Springer, 2005.
\bibitem{Munkres} J.~R.~Munkres, \emph{Elements of Algebraic Topology}, Addison-Wesley, 1984.
\bibitem{Quillen} D.~Quillen, Homotopy properties of the poset of nontrivial $p$-subgroups of a group, \emph{Adv.\ Math.} \textbf{28} (1978), 101--128. doi:10.1016/0001-8708(78)90058-0.
\bibitem{Shannon48} C.~E.~Shannon, A mathematical theory of communication, \emph{Bell System Tech.\ J.} \textbf{27} (1948), 379--423 and 623--656.
\bibitem{UNNSGIG} UNNS Research Collective, \emph{UNNS Generative Information Geometry}, research collection of 15 releases with README, repository ukbbi/UNNS, path projects/UNNS\_Generative\_Information\_Geometry, commit a115e4bc6b8ee734e6146b3b17a6b121e65d31c3.
\bibitem{UNNSFound} UNNS Research Collective, \emph{Mathematical Foundations of the UNNS Substrate}, reference collection v001, repository ukbbi/UNNS, path projects/UNNS\_MATH\_FOUNDATIONS, commit bb8b12d753c600c597feb00df17fe139c7cc65fa. \url{https://github.com/ukbbi/UNNS/tree/bb8b12d753c600c597feb00df17fe139c7cc65fa/projects/UNNS_MATH_FOUNDATIONS}
\bibitem{Ziegler} G.~M.~Ziegler, \emph{Lectures on Polytopes}, Graduate Texts in Mathematics \textbf{152}, Springer, 1995.
\end{thebibliography}

\end{document}
